\scanpage{9}
\mspage{\textcircled{2}}
% letterhead
\begin{center}
\deleted{ACADEMIE D'AIX-MARSEILLE}\\
\rule{4em}{0.4pt}\\
\deleted{INSPECTION ACADEMIQUE}\\
\deleted{DE VAUCLUSE}\\
\rule{1.6em}{0.4pt}
\end{center}
% ll. 1–5
$\rho$ définit $\rho^*:F^{\widetilde\Delta_m}\longrightarrow F^{\widetilde\Delta_n}$, et la relation $\rho_\pi^{\unsure{*|\prime}}(\beta)=\alpha$ signifie que $\rho^*\mid\beta$ se factorise par $\alpha$, i.e. $\rho^*$ induit
\[
 \unsure{\varphi|\Phi}(\rho):
 \underset{\substack{\rotatebox{90}{$=$}\\\varphi(\widetilde\Delta_m,\beta)}}{\beta}\longrightarrow
 \underset{\substack{\rotatebox{90}{$=$}\\\varphi(\widetilde\Delta_n,\alpha)}}{\alpha}.
\]
% ll. 6–10
On a ainsi défini une loi fonctorielle $\varphi$. Le foncteur $\varphi$ se prolonge par continuité en
\[
 \Phi:\underbrace{\bigl(\overbrace{(\widetilde\Delta_{/\pi})^{\mathrm o}}^{Z}\bigr)^\wedge}
 _{\substack{\rotatebox{90}{$=$}\\(\widetilde\Delta_{/\pi})^\vee=Z^\wedge}}
 \longrightarrow Y^\wedge,
\]
foncteur commutant aux $\varinjlim$ quelconques
% ll. 11–15
Je suspecte que \uline{ce foncteur est exact à gauche},$^*$ donc qu'il correspond à un morphisme de topos
% slanted
\margin{$^*$ c'est ici que pourrait intervenir l'hypothèse de \unsure{connexité|connexion} faite \unsure{sur|en} $X$ (et \unsure{laquelle est|qui est-elle} \unsure{requise : ?|requise ?}), mais pas encore}
\[
 f:Y^\wedge\longrightarrow Z^\wedge
\]
\[
 (\text{où}\qquad Z\overset{\text{déf}}{=}(\widetilde\Delta_{/\pi})^{\mathrm o}).
\]
% archive sheet number: 8
